\section{A positive secant kernel for every Euler remainder}
\label{sec:secant}
Let $a,b>0$, $w=a+b>0$, and let $R_N=2^N(E_N-g_{a,b})$
use the Euler partial sums defined in \eqref{eq:eulerdef}, now for arbitrary $a,b>0$. Define
\[
 H_N(s)=\frac{(1-s^{-2})^N}{1+s^{-2}},\qquad s\ge1.
\]
\begin{theorem}[Positive secant identity]
For every $a,b>0$, the Euler series converges absolutely to the analytic
value $g_{a,b}$. Its partial sums decrease strictly to that value.
For every integer $N\ge1$, with $V\sim\mathrm{Beta}(a,b)$,
\begin{equation}\label{secant:eq:remainder}
 R_N=\mathbb E\!\left[\frac1{\Gamma(w)}\int_0^\infty
 r^{w-1}\frac{H_N(e^r)-H_N(e^{Vr})}{e^r-e^{Vr}}\,dr\right].
\end{equation}
The quotient is positive for $0<V<1$. Its extension to $V=1$ is
$H_N'(e^r)$. Formula~\eqref{secant:eq:remainder} gives a positive
integral. When a finite signed Hausdorff representation exists, its
formula for the same remainder instead uses cancellation between
positive and negative parts. The positive secant formula also works
below the signed-measure threshold.
\end{theorem}
\begin{proof}

Start with the positive Mellin representation, where $f(n)=f_n$,
\[
 f(n)=\frac1{\Gamma(a)\Gamma(b)}\int_0^\infty\!\int_0^\infty
 u^{a-1}t^{b-1}\frac{e^{-u}}{e^t-1}
 \bigl(e^{-2nu}-e^{-2n(u+t)}\bigr)\,du\,dt.
\]
This is the finite geometric-sum formula for $H_{2n}^{(b)}$ followed
by the gamma integral for $(2n+1)^{-a}$. Applying the finite binomial
transform gives
\[
 -d_k=\frac1{\Gamma(a)\Gamma(b)}\int_0^\infty\!\int_0^\infty
 u^{a-1}t^{b-1}\frac{e^{-u}}{e^t-1}
 \left((1-e^{-2(u+t)})^k-(1-e^{-2u})^k\right)du\,dt.
\]
Every displayed integrand is nonnegative. Tonelli's theorem therefore
permits summing $\sum_{k\ge N}2^{N-k-1}(-d_k)$ under the double
integral. Summing the two geometric tails produces
\[
 \frac{e^{-u}}{e^t-1}
 \left(H_N(e^{u+t})-H_N(e^u)\right)
 =\frac{H_N(e^{u+t})-H_N(e^u)}{e^{u+t}-e^u}.
\]
The beta change of variables $r=u+t$, $v=u/r$ gives the asserted
integral for this nonnegative series. To establish its finiteness and
identify the Euler sum without a restriction on $w$, first take $N=1$.
Direct simplification shows that its secant kernel is twice the
kernel in \eqref{eq:Jw}. Therefore
\[
 \sum_{k\ge1}\frac{-d_k}{2^{k+1}}
   =\mathbb E J_w(V)=-g_{a,b}<\infty.
\]
This proves absolute convergence to the analytic value for every
positive pair. Also $d_k<0$ for $k\ge1$, since its displayed positive
integrand is strictly positive. Hence the partial sums decrease
strictly. Every later geometric tail is finite and equals
$2^N(E_N-g_{a,b})$, proving the formula for all $N\ge1$.
No conditionally convergent rearrangement is needed.
\end{proof}

For $N=1$, $H_1(s)=(s^2-1)/(s^2+1)$ and
\[
 H_1''(s)=\frac{4(1-3s^2)}{(1+s^2)^3}<0\qquad(s\ge1).
\]
Consequently its secant slope decreases as the lower endpoint $e^{vr}$
increases. Direct simplification gives
\[
 \frac{H_1(e^r)-H_1(e^{vr})}{e^r-e^{vr}}
 =\frac{e^{-vr}+e^{-r}}{2\cosh(vr)\cosh r},
\]
which is exactly twice the kernel used in the allocation theorem.
For $N=2$, however, $H_2(1+h)=2h^2+O(h^3)$, so $H_2$ is convex near
one. The same pointwise concavity proof does not extend to every $N$.

There is a particularly small exact obstruction to monotonicity in Euler
order away from the critical line. At $(a,b)=(1,1)$,
\[
 g_{1,1}=\operatorname{Im}\frac12\log^2(1-i)
       =-\frac{\pi\log2}{8},\qquad f(1)=\frac{H_2}{3}=\frac12.
\]
Since $E_1=0$ and $E_2=-f(1)/4=-1/8$,
\[
 R_1=\frac{\pi\log2}{4},\qquad
 R_2=\frac{\pi\log2}{2}-\frac12,
 \qquad R_2-R_1=\frac{\pi\log2-2}{4}>0.
\]
The last inequality follows already from $\pi>3$ and
$\log2=2\operatorname{artanh}(1/3)>2/3$.
This does not challenge the critical-line monotonicity theorem, whose
endpoint atom supplies a different argument. It shows why the sharp
bound for $-2g$ and the first Euler remainder cannot be declared the
sharp bound for every Euler order without an additional argument.
